% Part: incompleteness
% Chapter: incompleteness-provability
% Section: godels-paper

\documentclass[../../../include/open-logic-section]{subfiles}

\begin{document}

\olfileid{inc}{inp}{gop}

\olsection{Bandhingan karo Makalah Asli G\"odel}

Ana gunane nyawisake wektu kanggo maca makalah G\"odel taun 1931.
Pambukane nggambarake gagasan-gagasan kang lagi wae dirembug.
Sanajan mung maca sakeplasan, gampang dingerteni apa kang kelakon ing
saben tahapan: kawiwitan G\"odel njlentrehake sistem formal $P$
(sintaksis, aksioma, aturan bukti); banjur dheweke nemtokake fungsi
lan relasi rekursif primitif; banjur dheweke nuduhake yen $x B y$
rekursif primitif, lan ngandharake yen fungsi lan relasi rekursif
primitif direpresentasikake ing $\Th{P}$. Sabanjure dheweke mbuktekake
teorema ora jangkep kaya ing ndhuwur. Ing Bagean 3, dheweke nuduhake
yen pratelan kang ora bisa dibuktekake bisa dijupuk minangka ukara
ing basa aritmetika. Iki asal-usule lema~$\beta$, kang uga kita
gunakake kanggo ngolah urutan nalika nuduhake yen fungsi rekursif
bisa direpresentasikake ing $\Th{Q}$. G\"odel ora nganti misahake
himpunan aksioma minimal kang cukup, nanging saiki kita ngerti yen
$\Th{Q}$ wis cukup. Pungkasane, ing Bagean 4, dheweke nggambarake
garis gedhe bukti teorema ora jangkep kapindho.

\end{document}
